% ─────────────────────────────────────────────────────────────
% Abstract
% ─────────────────────────────────────────────────────────────
\begin{abstract}
\noindent
The discourse on artificial intelligence is dominated by two claims that are rarely separated: that AI development is hitting technical limits, and that leading laboratories are shaping regulation to protect their market position. This essay argues that neither claim, stated in that form, is analytically useful. The relevant change is not a ceiling on capability but a shift in the economic gradient of deployment: the cost of fixed capability is collapsing, production traffic is migrating to commoditized open-weight models, and the premium for frontier capability is decaying. Drawing on Carlota Perez's theory of technological revolutions, the essay interprets this shift as a transition from product competition toward infrastructure competition---from margin extracted through differentiation to rent extracted through control over the conditions of access. It then examines the institutional response: the Anthropic/METR embedded-evaluator arrangement and the draft bill that would generalize it. The essay does not infer motive. It evaluates design. An eight-dimension framework distinguishes observed design features from the mechanisms that may produce exclusionary effects and from the evidence that would demonstrate such effects. The framework is then validated against three further cases: the EU AI Act's conformity-assessment regime, China's algorithm registration regime, and the UK AI Safety Institute. Applied to the Anthropic template, the framework finds the arrangement contestable on publication and potentially exclusionary on access, proportionality, and open-weight treatment, while acknowledging that some design features may be justified by genuine safety requirements. The historical cases---railways, the internet, electrification, nuclear power---do not predict capture; they show that distributional outcomes depend on design choices being made now. The contribution is a mechanism and an evaluative framework, not a verdict.
\end{abstract}

\bigskip
\noindent\textbf{A note on evidence.} The empirical claims in this essay draw on industry production data (Vercel's AI Gateway Index and Epoch AI's inference-cost analysis) rather than peer-reviewed sources. The traffic data derive from a single gateway provider and should be treated as a convergent indicator, not a market-wide measure. Where industry analyses are used, their limitations are stated. Where primary documents are cited (Anthropic, METR, legislative drafts), URLs and access dates are provided in the reference list.

\medskip
\noindent\textbf{Keywords:} AI scaling; Carlota Perez; techno-economic paradigms; infrastructure rent; regulatory design; incumbent participation; open-weight models; evaluation regimes; EU AI Act; algorithmic governance.
